\section{Econometric Framework}
\label{sec:framework}
\subsection{Structural moments and orthogonalization}
Let $Y_i$ be the outcome, $D_i$ endogenous treatment, $X_i$ controls, and $Z_i\in\mathbb R^{d_z}$ instruments. Under the structural-rank assumptions of \textcite{chernozhukov_hansen_2005}, the monotone structural quantile function has the representation
\begin{equation}
    Y_i = q(U_i,D_i,X_i),
    \qquad
    U_i \mid X_i,Z_i \sim \operatorname{Uniform}(0,1),
    \label{eq:structural-ivqr}
\end{equation}

\begin{equation}
    q(\tau,D_i,X_i)
    = c_0(\tau)
      + D_i^{\prime}\alpha_0(\tau)
      + X_i^{\prime}\beta_0(\tau).
    \label{eq:linear-structural-quantile}
\end{equation}
The treatment coefficient $\alpha_0(\tau)$ is the target. The intercept and control coefficients are nuisances. Rank similarity permits treatment-state-specific ranks with the same relevant conditional distribution; rank invariance is the stronger common-rank case. The structural quantile at fixed treatment generally differs from the observational quantile conditional on endogenous treatment. Monotonicity and conditional rank uniformity imply
\begin{equation}
    \operatorname{E}\!\left[
        \tau
        - \mathbf{1}\!\left\{
            Y_i \leq c_0(\tau)
            + D_i^{\prime}\alpha_0(\tau)
            + X_i^{\prime}\beta_0(\tau)
        \right\}
        \mathrel{\big|} X_i,Z_i
    \right]
    = 0.
    \label{eq:ivqr-conditional-moment}
\end{equation}
Instrument validity supports this restriction; identification additionally requires variation that distinguishes alternative treatment coefficients.

For candidate $a$, define $\widetilde X_i=(1,X_i')'$, $\rho_\tau(v)=v[\tau-\mathbf1\{v<0\}]$, and the population nuisance profile
\begin{equation}
    \begin{pmatrix}
        c(a)\\
        \beta(a)
    \end{pmatrix}
    = \operatorname*{arg\,min}_{c,b}
      \operatorname{E}\!\left[
          \rho_{\tau}\!\left(Y_i-D_i^{\prime}a-c-X_i^{\prime}b\right)
      \right].
    \label{eq:profile-beta}
\end{equation}

\begin{equation}
    \epsilon_i(a)
    =Y_i-D_i^{\prime}a-c(a)-X_i^{\prime}\beta(a).
    \label{eq:profile-residual}
\end{equation}

\begin{equation}
    \begin{aligned}
        M(a)
        &=\operatorname{E}\!\left[
            Z_i\widetilde X_i^{\prime}f_{\epsilon(a)}(0\mid X_i,Z_i)
        \right],
        & M(a)&\in\mathbb{R}^{d_z\times(p+1)},\\
        J(a)
        &=\operatorname{E}\!\left[
            \widetilde X_i\widetilde X_i^{\prime}f_{\epsilon(a)}(0\mid X_i,Z_i)
        \right],
        & J(a)&\in\mathbb{R}^{(p+1)\times(p+1)}.
    \end{aligned}
    \label{eq:profile-MJ}
\end{equation}

\begin{equation}
    \delta(a)=M(a)J(a)^{-1},
    \qquad
    \delta(a)\in\mathbb{R}^{d_z\times(p+1)},
    \label{eq:profile-delta}
\end{equation}

\begin{equation}
    \Psi_i(a)
    =Z_i-\delta(a)\widetilde X_i
    =Z_i-d_0(a)-\Gamma(a)X_i.
    \label{eq:profile-residualized-instrument}
\end{equation}
Here $f_{\epsilon(a)}(0\mid X_i,Z_i)$ is the conditional residual density at zero. With $\eta=(c,\beta',\operatorname{vec}(\delta)')'$, let $\eta(a)$ collect the candidate-specific profile and $\eta_0=\eta(\alpha_0(\tau))$. The theoretical score is
\begin{equation}
    \begin{aligned}
        g_{\tau}(W_i;a,\eta)
        &=
        \left[
            \tau-\mathbf{1}\!\left\{
                Y_i\leq D_i^{\prime}a+c+X_i^{\prime}\beta
            \right\}
        \right]
        \left(Z_i-\delta\widetilde X_i\right),\\
        g_{\tau}(W_i;a,\eta(a))
        &=
        \left[
            \tau-\mathbf{1}\!\left\{
                Y_i\leq D_i^{\prime}a+c(a)+X_i^{\prime}\beta(a)
            \right\}
        \right]\Psi_i(a),
    \end{aligned}
    \label{eq:orthogonal-score}
\end{equation}

\begin{equation}
    \operatorname{E}\!\left[
        g_{\tau}\!\left(W_i;\alpha_0(\tau),\eta_0\right)
    \right]
    =0.
    \label{eq:orthogonal-score-moment}
\end{equation}

\begin{equation}
    \left.
    \partial_{\eta}
    \operatorname{E}\!\left[
        g_{\tau}\!\left(W_i;\alpha_0(\tau),\eta\right)
    \right]
    \right|_{\eta=\eta_0}
    =0.
    \label{eq:neyman-orthogonality}
\end{equation}
Under finite density-weighted moments, a nonsingular $J$, and the requisite differentiation conditions, the derivative in the intercept and slopes is $-[M_0-\delta_0J_0]=0$. The residualization derivative vanishes by the quantile moment. Appendix~\ref{app:score-algebra} gives the algebra. These are properties of the ideal construction under its conditions, rather than guarantees for every executable implementation.

This distinction is material for the present DGP. At truth, the outcome residual is $e_\tau=D(u-q_\tau)$, with $q_\tau=\Phi^{-1}(\tau)$ and $D=\Phi(T)$, $T=d^*(\kappa)$. Conditional on $u=q_\tau$, $X$, and observed $Z$, the variance of the normal index is $2.91-\kappa^2>1$ for every investigated setting. The density at zero contains a term proportional to $\operatorname E[1/\Phi(T)\mid u=q_\tau,X,Z]$, which is not finite. The experiment therefore does not establish the ideal density-weight regularity conditions or validate the associated asymptotic theory.

\subsection{Preserved estimators and inference}
\label{sec:implementation}
Oracle-GMM uses $X_1,\ldots,X_{10}$, the designated oracle set. Full-GMM uses all 100 controls. Both fit an unpenalized candidate-specific quantile regression with an intercept, then directly invert density-proxy-weighted control moment matrices without a residualization intercept. DML-IVQR fits quantile Lasso and regularized instrument regressions using all controls. Its Belloni--Chernozhukov pivotal penalty follows the executable route of \textcite{chen_huang_tien_2021}; their comparison with cross-validation motivates retaining this route.

The quantile-Lasso objective is an unscaled sum of check losses plus one half of the supplied coefficient-specific penalty vector times absolute coefficients, including the intercept. There is no post-Lasso refit. Instrument regressions use square-root-density-transformed instruments and controls with an ordinary unpenalized intercept. The score subtracts the retained intercept and original-control fitted slopes from the untransformed instrument. This intercept differs from the constant regressor in the ideal weighted projection. The fitted-normal proxy is evaluated at candidate residuals; the preserved R call \texttt{dnorm(e, mean(e), var(e))} supplies the residual variance where R interprets a standard deviation. It is not a direct estimate of the theoretical conditional density at zero.

Nuisances and scores use the same observations, with no sample splitting or cross-fitting. Fresh pivotal draws are generated at each penalty call, introducing algorithmic randomness across candidate evaluations. Appendix~\ref{app:reproducibility} records the penalty formulas, residualization rules, and environment details. The estimator comparison encompasses these implementation differences as well as regularization and nuisance dimension; it cannot isolate the effect of regularization alone.

For estimator $e$, let $\widehat g_{e,i}(a)$ be its score. The common criterion uses the uncentered score second moment:
\begin{equation}
    \widehat g_{e,n}(a)
    =
    \frac{1}{n}\sum_{i=1}^{n}\widehat g_{e,i}(a).
    \label{eq:sample-gmm-moment}
\end{equation}

\begin{equation}
    \widehat\Sigma_e(a)
    =
    \frac{1}{n}\sum_{i=1}^{n}
    \widehat g_{e,i}(a)\widehat g_{e,i}(a)^{\prime}.
    \label{eq:score-covariance}
\end{equation}

\begin{equation}
    W_{e,n}(a)
    =
    n\widehat g_{e,n}(a)^{\prime}
    \widehat\Sigma_e(a)^{-1}
    \widehat g_{e,n}(a).
    \label{eq:gmm-criterion}
\end{equation}

\begin{equation}
    \widehat\alpha_e(\tau)
    =
    \operatorname*{arg\,min}_{a\in\mathcal A}W_{e,n}(a).
    \label{eq:gmm-point-estimator}
\end{equation}
The Monte Carlo candidate set $\mathcal A$ is the fixed grid specified below; the first exact minimizer is retained, and an incomplete profile receives no point estimate. Direct inversions use neither a generalized inverse nor additional matrix regularization. Failures or non-finite statistics are recorded as numerical failures.

The procedures apply the asymptotic $\chi^2_{d_z}$ reference cutoff to each candidate statistic. At significance level $\gamma$, inversion gives
\begin{equation}
    \mathcal C_{e,1-\gamma}(\tau)
    =
    \left\{
        a\in\mathcal A:
        W_{e,n}(a)\leq c_{1-\gamma,d_z}
    \right\}.
    \label{eq:robust-confidence-region}
\end{equation}
Here $c_{1-\gamma,d_z}$ is the chi-square quantile, with $d_z=2$. The corresponding theoretical candidate tests can accommodate weak identification under their regularity conditions \parencite{chernozhukov_hansen_2008}; finite-sample calibration of these preserved implementations remains an empirical question. This construction differs from exact conditional finite-sample inference \parencite{chernozhukov_hansen_jansson_2009}. Weak-identification robustness does not ensure narrow sets or high power, and a finite numerical domain cannot establish global boundedness.
